\chapter*{第二部分收束：从数据处理到可建模样本}
\addcontentsline{toc}{chapter}{第二部分收束：从数据处理到可建模样本}

完成 Part II 后，我们已经接触了五类天文与物理数据：带误差和单位的测量值、星表、图像、光谱、时间序列，以及带模型假设的实验或模拟数据。它们看起来各不相同，但如果要进入后面的机器学习和案例项目，都必须先完成同一个转化：从“原始或教学数据”变成“可解释、可检查、可建模的样本”。

这一步不能跳过。模型只能看见你交给它的输入；如果输入的物理含义、误差来源、采样边界和选择效应没有被整理清楚，后面再复杂的算法也只能在混乱基础上给出更复杂的混乱。

\section*{一条通用的数据处理算法}

你可以把 Part II 的共同工作流写成下面这条算法：

\begin{enumerate}
    \item \textbf{定义科学问题}：先说清楚要比较、估计、分类或检验什么；
    \item \textbf{写下数据契约}：列出必需字段、单位、坐标、采样方式和元数据；
    \item \textbf{读入并验证数据}：检查缺失值、非法值、数组形状、header 和基本统计；
    \item \textbf{建立误差预算}：至少区分随机误差、系统误差和由处理选择引入的不确定性；
    \item \textbf{执行类型相关处理}：例如星表派生列、图像背景扣除、光谱归一化、时间序列折叠或实验模型拟合；
    \item \textbf{用图表诊断结果}：主图展示结构，残差图或质量图检查异常；
    \item \textbf{记录选择边界}：写下筛选条件、参数选择、失败样本和暂未处理的偏差；
    \item \textbf{产出可建模样本}：把最终进入模型或报告的字段、标签、图表和运行入口记录清楚。
\end{enumerate}

这条算法比任何单个函数都重要。因为真实项目中函数会换、库会升级、数据会变大，但这八步所保护的科学边界不会消失。

\section*{不同数据类型的最小交付件}

\begin{center}
\begin{tabular}{p{0.18\textwidth} p{0.34\textwidth} p{0.34\textwidth}}
\toprule
数据类型 & 最小处理结果 & 进入建模前最该检查什么 \\
\midrule
星表 & 清洗后的对象表、派生列、筛选说明 & 字段单位、视差或距离质量、选择效应 \\
图像 & 背景修正信号、源位置、cutout 或测光表 & WCS/header、背景估计、孔径和像素尺度 \\
光谱 & 归一化光谱、谱线位置、线强或等效宽度 & 波长校准、采样一致性、信噪比和连续谱处理 \\
时间序列 & 候选周期、相位折叠曲线、残差或周期排名 & cadence、窗口函数、混叠峰和候选稳定性 \\
实验/模拟 & 参数估计、残差图、卡方或 Monte Carlo 分布 & 模型假设、误差模型、系统偏差和外推边界 \\
\bottomrule
\end{tabular}
\end{center}

这张表也可以当作课程项目中的检查表：每当你准备把某类数据交给机器学习模型之前，先问自己是否已经产出了相应的最小处理结果，并检查了最容易破坏解释的环节。

\section*{从数据处理到机器学习的桥}

Part III 会引入训练集、验证集、特征、标签、损失函数和模型评估。进入那里之前，你需要带走 Part II 的三个结论。

第一，\textbf{特征不是中性的列}。一个特征可能来自直接观测，也可能来自公式推导、背景扣除、归一化、重采样或模型拟合。不同来源的特征带着不同误差和偏差。

第二，\textbf{标签不是天然真理}。星表类别、光谱类型、周期类型或实验状态都可能来自人工判断、阈值规则、已有目录或模拟设定。标签来源会影响模型能学到什么，也会影响模型错误应该怎样解释。

第三，\textbf{训练/测试划分不能修复数据理解错误}。如果样本筛选本身已经有强选择效应，或者图像和光谱在预处理阶段泄漏了目标信息，后面的验证分数可能仍然很好看，但科学意义已经不可靠。

\begin{warnbox}[不要把 Part II 当成模型前处理的附录]
数据处理不是机器学习之前的机械清洗，而是科学问题进入模型之前的语义整理。只有当字段、单位、误差、坐标、采样和处理记录都足够清楚时，模型评估才有资格被解释成科学证据。
\end{warnbox}

\section*{进入 Part III 前的自检}

在继续学习机器学习之前，建议你能完成下面的检查：

\begin{itemize}
    \item 对任意一个表格样本，写出字段契约和至少一个派生量的公式来源；
    \item 对任意一个图像 cutout，说明背景、孔径、像素坐标和天空坐标分别是什么；
    \item 对任意一条光谱，解释归一化、谱线位置和信噪比如何影响物理判断；
    \item 对任意一组时间序列，说明采样窗口为什么可能制造假周期；
    \item 对任意一次实验拟合，区分模型参数、残差结构和不确定度估计；
    \item 对一个准备进入机器学习的样本表，写出哪些列是原始观测量，哪些列是派生特征，哪些列可能泄漏答案。
\end{itemize}

如果这些问题能够回答，Part III 的机器学习就不再只是“调用模型”。它会变成一个更清楚的问题：在已经理解数据来源和物理边界的前提下，模型能帮助我们发现什么结构、做出什么预测，以及哪些地方仍然不能相信。
